% Generated from validated Article Draft Result. Do not hand-edit; revise the structured draft instead.
\section{Reasoningが標準機能になる — 「考えるモデル」から実行能力の一層へ}
\noindent\textbf{2025年前半、推論能力は単独のモデル分類から、API・ツール利用・長文脈・視覚理解と結び付く実行能力へ移り始めた。重要なのは順位表ではなく、どのsurfaceで、どの形のthinkingが利用可能になったかである。}\autocite{src-421ee9af3f80,src-650b567698aa,src-8f8555b73e1b}

1月20日のDeepSeek-R1 API提供は、推論モデルを専用endpointで利用するservice eventとして現れた。一方、3月24日のDeepSeek-V3-0324はdeepseek-chat系の更新であり、R1の登場とV3系service lifecycleは別の出来事として扱う必要がある。\autocite{src-8f8555b73e1b}


1月末のo3-miniはChatGPTとAPIの双方へ小型reasoning modelを持ち込み、4月のo3/o4-miniではreasoningとtool accessの結合がさらに明示された。ここではモデル名の更新以上に、推論がdeveloper-facing executionへ接続されたことが重要である。\autocite{src-650b567698aa,src-bf92f3d905c3}


2月24日のClaude 3.7 Sonnetはhybrid reasoningを掲げ、同じ発表でClaude Codeのresearch previewも提示された。モデル内部のthinkingと、terminal上で作業するcoding agentは同一物ではないが、同じ日に接続されたことがH1の転換を象徴する。\autocite{src-a6c92f94a0af}


3月25日のGemini 2.5 Pro Experimentalはthinking、long context、coding、multimodal capabilityを一つのexperimental modelとして提示した。H1ではこのexperimental availability自体が重要で、6月のGAは後段のlifecycleとして別に保持する。\autocite{src-421ee9af3f80}


QwenのQVQ-Maxは3月末にvisual reasoningの系譜を示した。reasoningという語がtext-only chainに閉じず、視覚入力へ広がっていた点は、同時期のnative multimodalityと合わせて読む必要がある。\autocite{src-8cb736b33646}


半年を通してみると、reasoningは独立した製品ジャンルというより、tool use、coding、research、multimodal processingを支えるcapability layerへ再分類できる。R1、o3系、Claude 3.7、Gemini 2.5は実装・提供形態が異なるため、単一benchmark順位だけで同列比較しない。\autocite{src-421ee9af3f80,src-8f8555b73e1b,src-a6c92f94a0af,src-bf92f3d905c3}


\begin{claimboundary}
本章で確定するのは、一次資料で確認できるrelease、availability、tool接続、model identityである。各社が示す性能優位やbenchmark評価は、独立再現された測定値としては扱わない。
\end{claimboundary}
